% Modèle de lettre de motivation une page (pdfLaTeX).
% 250 à 350 mots. Remplacer tout ce qui est entre [crochets].
\documentclass[11pt,a4paper]{article}
\usepackage[utf8]{inputenc}
\usepackage[T1]{fontenc}
\input{glyphtounicode}
\pdfgentounicode=1
\usepackage[scaled=0.95]{helvet}
\renewcommand{\familydefault}{\sfdefault}
\usepackage[left=1in, right=1in, top=0.9in, bottom=0.9in]{geometry}
\usepackage{xcolor}
\usepackage{hyperref}
\usepackage{microtype}

\definecolor{linkblue}{RGB}{70, 120, 142}
\hypersetup{colorlinks=true, urlcolor=linkblue}

\setlength{\parindent}{0pt}
\setlength{\parskip}{7pt}
\linespread{1.12}
\hyphenpenalty=10000
\exhyphenpenalty=10000
\emergencystretch=3em

\begin{document}
\pagestyle{empty}

% --- EXPÉDITEUR ---
[Prénom Nom]\\[4pt]
[Ville], France\\[4pt]
[Téléphone]\\[4pt]
\href{mailto:[email]}{\underline{[email]}}

\vspace{20pt}

% --- OBJET ---
{\raggedright \textbf{Objet~: Candidature, Stage [année], [Intitulé exact], [Entreprise], [Lieu]}\par}

Madame, Monsieur,

% 1. Accroche : un fait précis et vérifié sur le site, le produit ou l'enjeu du poste, relié à la mission.
[Accroche, 2 à 3 phrases.]

% 2. Preuves : 2 ou 3 exigences de l'offre, chacune prouvée par une expérience réelle.
[Exigence 1 → expérience : contexte, action, résultat.]

[Exigence 2 → expérience : contexte, action, résultat.]

% 3. Ce qui manque, dit sobrement.
[Je n'ai encore jamais fait X~: c'est ce que je viens apprendre.]

% 4. Centre d'intérêt relié à une exigence du poste.
[Une ou deux phrases.]

% 5. Clôture : date de début, durée, écart éventuel avec l'offre.
Mon stage de fin d'études peut commencer en [mois année], pour [durée] mois. Je serais heureux d'en échanger avec vous lors d'un entretien.

Je vous prie d'agréer, Madame, Monsieur, l'expression de mes salutations distinguées.

[Prénom Nom]

\end{document}
